% !TeX program = xelatex
\documentclass[12pt,a4paper]{ctexart}

\usepackage{amsmath,amssymb,amsfonts}
\usepackage{graphicx}
\usepackage{booktabs}
\usepackage{geometry}
\usepackage{float}
\usepackage{caption}
\usepackage{subcaption}
\usepackage{hyperref}
\usepackage{enumitem}
\usepackage{array}
\usepackage{longtable}
\usepackage{multirow}
\usepackage{setspace}
\usepackage{ragged2e}
\usepackage{pdflscape}

\geometry{left=2.2cm,right=2.2cm,top=2.2cm,bottom=2.2cm}
\setlength{\parskip}{0.35em}
\onehalfspacing
\hypersetup{colorlinks=true,linkcolor=blue,citecolor=blue,urlcolor=blue}
\emergencystretch=2em
\newcommand{\imginclude}[2][]{\includegraphics[#1]{results/#2}}
\newcommand{\modeltrio}{baseline 与 fear-memory}
\newcommand{\trainrmse}{\mathrm{RMSE}_{\mathrm{train}}}
\newcommand{\valrmse}{\mathrm{RMSE}_{\mathrm{val}}}

\title{\textbf{恐惧效应下捕食者--猎物动力学：\\机制路径、预测比较与参数可辨识性局限}}
\author{%
  贺小轩（PB23151820）\quad 李松茂（PB23151824）\\[0.25em]
  王玉麟（PB23151808）%
}
\date{}

\begin{document}
\maketitle

\begin{abstract}
经典捕食者--猎物模型通常把捕食风险等同于直接消耗，但猎物还会因感知风险改变繁殖、觅食和活动。本文以逻辑斯蒂--Holling II 基线与带指数记忆的繁殖抑制模型为核心。为减少默认参数强度差异的干扰，五种恐惧机制均相对无恐惧基线，在参考正平衡点处校准为 20\% 标准化抑制。真实数据比较采用前 80\% 拟合、后 20\% 连续多步留后验证，并统一使用归一化 RMSE、AICc、拟合状态过滤和参数 profile 诊断。

结果表明，标准化基准下不同恐惧机制产生不同的长期种群响应，结论具有明显的机制路径依赖性。记忆模型的 $\phi$ 扫描 31/31 在当前积分范围内未灭绝，其中 11 个情景尚未通过长期收敛诊断。15 个拟合单位（序列级）中，含恐惧模型在留后预测上各有胜负；在 baseline 与 fear-memory 的两模型比较中，baseline 在 13 条序列上获 AICc 支持，扩展到六模型后则在 8 条序列上取得最低 AICc。全序列 profile 分析同时显示记忆相关参数普遍存在宽可接受范围。因此，复杂模型可以改善部分序列的预测，但预测改善不能直接解释为恐惧机制真实存在，也不能支持对真实恐惧强度的稳定估计。

\textbf{关键词：}捕食者--猎物；恐惧效应；记忆效应；留后验证；可辨识性
\end{abstract}

\tableofcontents
\newpage

\section{引言}

捕食者不仅通过捕食直接减少猎物数量，也会通过风险线索改变猎物繁殖、觅食、活动和栖息地选择。这类非消耗性效应（non-consumptive effect, NCE）可能进一步改变捕食者资源、种群共存和系统稳定性。问题的难点不在于是否加入一个“恐惧参数”，而在于恐惧以何种机制进入方程，以及数据是否足以区分这些机制。

本文围绕三个问题展开。第一，不同恐惧机制是否产生一致的种群动力学结论？第二，含恐惧模型能否改善真实时间序列的留后预测？第三，当模型复杂度增加时，数据能否可靠识别恐惧参数？

为回答这些问题，本文首先建立 Holling II 基线与 fear-memory 主模型，分析恐惧强度、记忆衰减和长期共存；随后在无恐惧基线的正平衡点处进行 20\% 标准化抑制校准，比较五种 Holling II 恐惧机制路径；最后对 15 个拟合单位（序列级）\footnote{含研究内重复：Andrén 七区、Killifish 三站和 Windermere 两湖盆分别来自同一研究，不能视为 15 个独立生态证据。详见附录 A。}进行六模型预测比较，并通过参数 profile 诊断检验预测改善是否真正意味着恐惧参数被识别。

\section{相关工作}

捕食者--猎物动力学的理论建模始于 Lotka--Volterra 方程，后经逻辑斯蒂增长和 Holling II 功能反应纳入环境容量与饱和捕食，成为该领域的标准框架。然而，这些经典模型仅通过直接捕食消耗连接两个种群，未考虑猎物对风险的感知与行为响应。

Abrams（2000）指出，适应性行为进入种群动力学方程的具体方式会改变预测的定性符号（如稳定性反转或共存状态变化），但该框架的结论高度依赖于模型参数化形式，未提供在真实数据中区分不同行为机制的操作性方法。Zanette 等（2003）通过经典鸣禽实验确证了感知捕食风险即可抑制繁殖产出，但其效应量是在受控实验中测得的，能否在宏观种群时间序列中被检测到仍属未知。Preisser 等（2007）的综合分析显示非消耗性效应的平均效应量与消耗性效应相当，但该比较基于实验操作的标准化效应量（Hedges' $d$），并非从种群动态观测中估计的连续恐惧强度参数。Wang 等（2019）引入了指数记忆核来参数化恐惧效应的持续性，并讨论了恐惧因子对稳定性的影响，但该工作的参数设定为 a priori 赋值，未进行真实数据校准或模型间交叉验证。Peacor \& Werner（2001）提出了性状介导间接效应（TMIE）的概念框架，Peacor 等（2022）进一步整理了相关实验证据，但这些证据本质上是对研究类型的计数清单（TMIE/NCE），并非可直接输入动力学模型的恐惧强度估计量。

这些工作共同提示两项关键的方法要求。其一，不同恐惧机制不能在各自默认参数下直接比较——参数选择的任意性会掩盖机制结构的真实差异；应相对无恐惧基线设置标准化抑制基准。该基准可减少通道内部强度差异，却不能使繁殖抑制与捕食率抑制等不同过程量完全可通约。其二，训练误差的降低不是机制识别的充分证据——复杂模型总是可以通过额外的自由度改善拟合；真实数据比较还需留后预测、复杂度惩罚（AICc）、优化状态过滤和参数可辨识性诊断。本文据此区分”模型能否预测轨迹”与”数据能否识别恐惧机制”。

\section{预备知识}

\subsection{状态、功能反应与恐惧抑制}

经典体系中，$x(t)$ 和 $y(t)$ 分别表示猎物与捕食者；记忆状态 $M(t)$ 表示过去捕食压力的指数加权累积。Holling II 捕食项为
\[
  H(x,y)=\frac{axy}{1+\theta x}.
\]

\subsection{模型比较指标}

两类观测分别使用训练段最大值确定的尺度归一化。总归一化 RMSE 为两类归一化残差合并后的均方根。训练段用于估计参数，末尾 20\% 数据仅用于从初始状态连续积分得到的多步留后预测。AICc 使用同一归一化训练残差，并按实际拟合参数数量惩罚复杂度：
\[
 \mathrm{AICc}=\mathrm{AIC}+
 \frac{2p(p+1)}{n-p-1}.
\]
其中 $p$ 为实际拟合参数数目，$n=2n_{\mathrm{train}}$ 为猎物与捕食者训练残差的总数。AIC 中的 $2p$ 是基础复杂度惩罚，AICc 最后一项进一步修正小样本下的乐观偏差；新增参数是否值得由完整的 $\Delta\mathrm{AICc}$ 判断，即拟合误差下降是否足以抵消两部分复杂度惩罚，而不能只看最后一项。该计算将归一化残差近似视为同方差 Gaussian 残差；生态时间序列中的自相关、过程噪声和观测误差未被显式建模，因此 AICc 仅作为当前候选模型集合内的相对比较指标。单次末尾 20\% 留后同样可能对切分点敏感。
优化状态分为成功、达到迭代上限但结果可用、失败。失败拟合保留用于诊断，但不进入模型胜负统计。

\section{模型与方法}
\label{sec:model-method}

\subsection{核心模型：Holling II baseline 与 fear-memory}

\paragraph{逻辑斯蒂--Holling II 基线}
\begin{equation}
\begin{aligned}
 \frac{dx}{dt}&=rx\left(1-\frac{x}{K}\right)-\frac{axy}{1+\theta x},\\
 \frac{dy}{dt}&=e\frac{axy}{1+\theta x}-\mu y.
\end{aligned}
\label{eq:baseline}
\end{equation}

\paragraph{繁殖抑制与指数记忆}
\begin{equation}
\begin{aligned}
 \frac{dx}{dt}&=rx\left(1-\frac{x}{K}\right)-\frac{axy}{1+\theta x}-r\phi Mx,\\
 \frac{dy}{dt}&=e\frac{axy}{1+\theta x}-\mu y,\\
 \frac{dM}{dt}&=y-\delta M.
\end{aligned}
\label{eq:memory}
\end{equation}
指数核通过线性链转化为状态 $M$；稳态时 $M^*=y^*/\delta$。$\delta$ 越大表示记忆衰减越快。

\subsection{扩展恐惧机制}

除核心 fear-memory 模型外，本文还引入四种 Holling II 框架下的替代恐惧机制以检验机制路径依赖性：
\begin{align}
 \text{fear-instant:}\quad
 \frac{dx}{dt}&=rx\left(1-\frac{x}{K}\right)-\frac{axy}{1+\theta x}-r\phi yx, \label{eq:instant}\\
 \text{fear-saturating:}\quad
 \frac{dx}{dt}&=r\left(1-\frac{\phi y}{y+h}\right)x\left(1-\frac{x}{K}\right)-\frac{axy}{1+\theta x}, \label{eq:saturating}\\
 \text{fear-foraging:}\quad
 \frac{dx}{dt}&=rx\left(1-\frac{x}{K}\right)-\frac{a x y}{(1+\psi y)(1+\theta x)}, \label{eq:foraging}\\
 \text{fear-handling:}\quad
 \frac{dx}{dt}&=rx\left(1-\frac{x}{K}\right)-\frac{a x y}{1+\theta(1+\psi y)x}. \label{eq:handling}
\end{align}
四者的捕食者方程均为相应捕食项乘以 $e$ 后减去 $\mu y$。fear-instant 表示瞬时繁殖抑制，fear-saturating 表示饱和繁殖抑制，fear-foraging 通过降低有效攻击率抑制觅食，fear-handling 则延长有效处理时间。

在理论标准化实验中，fear-memory 与 fear-instant 使用相同的 $\phi=0.00783431$ 和 $\delta=1$：在正平衡点处，由于 $M^*=y^*/\delta=y^*$，恐惧项 $r\phi M x$ 退化为 $r\phi y x$，与 fear-instant 完全等价\footnote{两者在平衡点附近产生不同的一阶修正：fear-memory 的指数记忆低通滤波 $dM/dt = y-\delta M$ 引入由 $\delta=1$ 定义的特征记忆时间尺度，离开平衡点后时变恐惧强度不同。}。真实数据拟合时，两模型的 $\phi$ 分别自由优化；因此其比较检验的是两种结构在相同参数预算下的预测表现，而不是固定相同 $\phi$ 后的纯记忆效应实验。

\subsection{五种机制与标准化基准比较}

所有恐惧机制在无恐惧正平衡点处校准为 20\% 抑制：繁殖路径匹配增长项抑制，觅食与处理时间路径匹配捕食率抑制。20\% 是人为选择的机制比较基准，目的是使不同通道造成的动力学差异清晰可见；它不代表文献共识，也不代表真实系统中的典型恐惧强度。由于不同路径匹配的是不同过程量，该设计是标准化基准比较，而不是严格相同生物学效应量的比较。校准参数与规则见附录 B。

\subsection{理论、数值与诊断方法}

ODE 使用 \texttt{solve\_ivp} 的 RK45 积分，并线性插值到均匀网格。长期指标来自末尾 20\% 时间窗；若相邻尾窗的均值、振幅或漂移相对变化超过 1\%，则延长积分，最多延长三次。灭绝判定使用相对系统尺度的阈值，并要求尾段持续低于阈值且无明确恢复。

\subsection{数据拟合与证据过滤}

15 条主序列包含哺乳类、鱼类与浮游动物，详见附录 A。正式基础比较逐序列拟合 baseline 与 fear-memory；完整扩展比较再加入 fear-instant、fear-saturating、fear-foraging 与 fear-handling。训练段为前 80\%，验证段为后 20\%；验证预测不是从验证起点重新初始化，而是从整条序列初始状态连续积分。比较顺序为：
\begin{enumerate}[label=(\arabic*),leftmargin=2em]
 \item 过滤失败拟合；
 \item 比较训练归一化 RMSE；
 \item 比较连续多步留后验证 RMSE；
 \item 比较统一归一化残差下的 AICc；
 \item 检查参数触边。
\end{enumerate}

baseline 逐序列拟合 6 个动力学参数（$r,K,a,\theta,e,\mu$）。五个 Holling II 恐惧候选模型均在相同核心参数基础上增加 1 个主要恐惧参数，因此各拟合 7 个参数。fear-memory 固定 $\delta=1$，其新增拟合参数为 $\phi$；$\theta$ 则始终作为 Holling II 核心参数自由拟合。fear-saturating 的 $h$ 固定为对应序列训练段捕食者中位数。初始状态取猎物与捕食者首次同时为正的观测，不作为自由参数。对 fear-memory 进一步设 $M(0)=y(0)$，使初始记忆与首个可观测捕食压力一致，避免引入弱可辨识的初始记忆状态。固定 $\delta=1$ 提供统一的记忆时间尺度（对应记忆半衰期 $\ln 2\approx 0.69$ 个序列时间单位），防止在短序列中与 $\phi$ 相互补偿导致两个恐惧参数均不可靠。

优化策略采用固定种子（0, 1, 2）的多起点差分进化，每个种子独立运行至收敛或达上限，选取最低目标函数值的解，再以 L-BFGS-B 进行一次有界局部精修。由于五个 Holling II 恐惧模型在恐惧强度为零时均退化为 baseline，拟合还显式计算“正式 baseline 核心参数+零恐惧强度”候选，并保证最终训练目标不高于该候选；即使零恐惧候选胜出，AICc 仍按 7 个参数惩罚。拟合仅接受在观测与留后积分范围内非负、有限且可计算的轨迹。因各序列时间尺度和观测含义不同，拟合得到的 $e,\mu$ 点估计仅用于本序列的轨迹拟合，不直接解释为跨系统可比较的生态速率。

为检验上述记忆设定是否主导结果，本文还对 $M(0)$ 与 $\delta$ 进行补充 profile。在离散网格上固定目标参数并局部重新优化其余参数，再使用渐近 $\chi^2_1$ 阈值得到近似 95\% profile 集。全部 15 条序列的 $M(0)/y(0)$ 近似 95\% profile 集均延伸至扫描区间两端；$\delta$ profile 中 14/15 条至少延伸至一个扫描边界。这说明训练数据未能在已扫描范围内给出封闭的可接受区间，参数变化可被其余参数补偿；扫描边界只是网格截断点，不是精确的区间端点，也不表示参数对轨迹没有影响。作为敏感性分析，本文还在同一留后段上比较不同 $\delta$ 候选，并报告其中最小的验证 RMSE；相对固定 $\delta=1$，其改善中位数约为 1.2\%。由于留后段参与了 $\delta$ 的选择，它不再是独立验证集，所得改善可能偏乐观，因此不用于模型胜负统计。完整结果见附录 E。

\section{理论与动力学分析}

\subsection{记忆模型的平衡关系}

在正平衡点处，捕食者方程给出
\[
 \frac{eax^*}{1+\theta x^*}=\mu,
\]
结合 $M^*=y^*/\delta$ 可得
\[
 x^*=\frac{\mu}{ea-\mu\theta},
 \qquad
 y^*=
 \frac{r(1-x^*/K)}
 {a/(1+\theta x^*)+r\phi/\delta}.
\]
正平衡要求 $ea>\mu\theta$ 且 $0<x^*<K$。因此在正平衡存在时，$x^*$ 不直接依赖 $\phi$，而 $y^*$ 随 $\phi$ 增大而下降。这解释了通过收敛诊断的 $\phi$ 扫描情景中猎物长期均值近似不变、捕食者均值下降的现象；未通过诊断的情景不能用该平衡关系解释为已收敛长期状态。

\section{实验与讨论}

\subsection{记忆恐惧与参数扫描}

默认情景的时间序列说明，恐惧与记忆会改变瞬态轨迹和捕食者长期水平（时间序列对照图见附录\,\ref{fig:main-ts}）。

$\phi\in[0,0.06]$ 的 31 个扫描情景均未触发灭绝判定。20 个情景通过长期收敛诊断，$\phi\geq0.040$ 的 11 个情景在最大延长后仍未通过；因此这些点只能说明当前积分范围内未灭绝，不能当作已收敛长期平衡。通过诊断的范围内，猎物长期均值约为 20.02，捕食者均值随 $\phi$ 增大而下降（$\phi$ 扫描图见附录\,\ref{fig:phi-scan}）。

\subsection{机制路径依赖性}

在当前参考参数与 20\% 标准化基准下，五种恐惧机制均保持共存，但长期均值方向明显不同（机制比较柱状图见附录\,\ref{fig:mechanisms}）。三种繁殖抑制路径使捕食者均值下降约 16.7\%--18.0\%，猎物均值基本不变；觅食抑制与处理时间延长同时提高猎物和捕食者均值。因此，种群变化的定性方向取决于恐惧的机制路径。由于不同路径匹配的是不同过程量，该结果展示路径依赖性，但不构成跨机制效应量的严格排序。

\subsection{真实数据的预测比较}

模型比较分为两层递进问题：(a) fear-memory 是否优于无恐惧 baseline？(b) 五种 Holling II 恐惧通道中是否有某一通道一致最优？

表\,\ref{tab:wins} 汇总 baseline 与 fear-memory 的序列级结果：留后验证中 baseline 在 9 条序列上最优，fear-memory 在 6 条上最优；AICc 则在 13 条序列上支持 baseline，在 2 条上支持 fear-memory。留后预测与复杂度校正后的胜负明显分歧，因此不能把预测优势直接解释为恐惧机制获得识别。

\begin{table}[H]
 \centering
 \caption{两种标准下的核心模型胜负汇总（序列级）}
 \label{tab:wins}
 \begin{tabular}{lccc}
  \toprule
 标准 & baseline & fear-memory & 并列\\
  \midrule
 留后验证归一化 RMSE 最优 & 9 & 6 & 0\\
 AICc 最优 & 13 & 2 & 0\\
  \bottomrule
 \end{tabular}
\end{table}

按相同训练段归一化、20\% 连续多步留后和固定种子多起点优化口径扩展到全部六模型后，90 次拟合均达到 success 或 usable\_limit 状态并进入主比较，其中 33 次为 success，57 次为 usable\_limit。五种 Holling II 恐惧通道均执行显式零恐惧候选检查；fear-memory、fear-instant、fear-saturating、fear-foraging 与 fear-handling 分别有 6、7、5、6、3 条序列选择零恐惧候选。零恐惧候选胜出时继承对应 baseline 的优化状态，避免隐藏核心参数的优化不确定性。六模型在序列级留后最优分布高度分散，仍无单一恐惧通道在留后预测和 AICc 上一致胜出。完整六模型胜负表、热力图及研究级聚合分析见附录 D。

为检验胜负是否为微小数值波动所致，进一步计算每条序列的 baseline 与最佳恐惧模型之间的 $\Delta\mathrm{AICc}$。结果表明，6/15 条序列的 $|\Delta\mathrm{AICc}|<2$，胜负在实质上可忽略；仅 2 条序列的 $|\Delta\mathrm{AICc}|>7$，且均支持恐惧模型。其余序列虽可能存在 AICc 最优模型，但证据差异未达到强证据阈值。

研究级敏感性分析（取研究内序列的指标中位数，将 Andrén 七区、Killifish 三站、Windermere 两湖盆各聚合为 1 个研究级单位，共 6 个研究级比较单位）显示，留后验证最优模型分别为 baseline、fear-memory、fear-instant、fear-foraging 与 fear-handling，其中 fear-instant 在 2 个单位上最优，其余各在 1 个单位上最优；AICc 最优分布为 baseline 4、fear-memory 1、fear-instant 1。胜负模式仍不支持单一通道一致胜出。因此，序列非独立性不改变本文的主要定性结论。完整六模型胜负表及研究级聚合结果见附录 D。

作为优化状态敏感性诊断，若仅保留 success 拟合，Windermere South 没有任何模型可比较，其余 14 条序列中每条保留的候选模型集合也不同。因此 success-only 结果不能形成候选集合一致的六模型胜负统计；赢家仍分散于多种模型，未出现单一通道一致胜出。该诊断同时说明，主胜负计数仍依赖 usable\_limit 解，解释时需保留这一计算不确定性。

Windermere 南北两湖盆的拟合误差在 15 条序列中最高（留后验证 RMSE 最高达 3.04），其原因与狗鱼的刺网移除干预有关——当捕食者被人为压制而猎物先爆发后自然回落时，Holling II 框架的参数估计因未建模的外部干预而畸变。完整案例分析见第\,\ref{sec:discussion}\,节讨论。

\begin{figure}[H]
  \centering
  \imginclude[width=0.92\textwidth]{seven_model_real_fits/report_protocol_six_model_validation_heatmap.png}
  \caption{15 条正式序列的六模型连续多步留后验证 RMSE。颜色表示相对该序列最佳可用模型的误差比。}
  \label{fig:six-model-heatmap}
\end{figure}

全部拟合均通过 success 或 usable\_limit 状态进入比较；其中超过半数为 usable\_limit，因此“可比较”仅表示目标值有限且轨迹可计算，不等同于优化器已收敛。

\begin{table}[H]
 \centering
 \caption{证据边界：当前数据在各个维度上能够支持与不能支持的结论}
 \label{tab:prediction-mechanism}
 \small
 \begin{tabular}{p{0.16\textwidth}p{0.35\textwidth}p{0.35\textwidth}}
  \toprule
  维度 & 证据支持的结论 & 证据不支持的结论\\
  \midrule
  预测能力 & 六模型比较中，含恐惧模型在 12/15 条序列上取得严格最低留后验证 RMSE；胜出通道因序列而异 & 添加恐惧一定改善种群预测；某一种通道普遍优于其他通道\\
  \midrule
  机制识别 & 不同恐惧通道在数值实验中产生不同定性方向；$\Delta\mathrm{AICc}$ 分布提示部分比较（6/15）在实质上不区分模型 & 真实系统中存在与某一特定恐惧通道对应的机制；数据偏好可以直接解释为机制证据\\
  \midrule
  恐惧强度估计 & fear-memory 的 $\phi$ 在 1/15 条序列中触及正式搜索上界；扩边后虽得到内部解，但验证误差变差；全序列 $M(0)$ 与 $\delta$ profile 显示记忆相关参数存在宽可接受范围 & 将 fitted $\phi$ 直接解释为真实恐惧强度；在缺少独立机制观测时跨序列比较恐惧抑制程度\\
  \midrule
  跨系统比较 & 6 个研究级聚合单位的胜负模式与序列级定性一致；哺乳类、鱼类和浮游动物均存在留后预测改善的个例 & 任一恐惧通道在某一类群中一致胜出；15 个拟合单位可视为 15 个独立生态证据\\
  \bottomrule
 \end{tabular}
\end{table}

\begin{figure}[H]
  \centering
  \imginclude[width=0.82\textwidth]{calibration_bda/figures/09_andren_lynx_roedeer_data_7_fear_memory.png}
  \caption{代表性拟合：Andrén 区 7 的 fear-memory 模型（AICc 最优）。阴影分界后的观测用于连续多步留后验证。}
  \label{fig:representative-fit}
\end{figure}

\subsection{恐惧记忆参数的可辨识性局限}
\label{sec:identifiability}

为检验 fear-memory 的恐惧与记忆参数是否被数据可靠识别，本文对 $\phi$、$M(0)$ 和 $\delta$ 进行补充诊断。$\phi$ 扩边诊断针对正式拟合中触及上界的序列；$M(0)$、$\delta$ 的 profile 覆盖全部 15 条序列。正文展示 3 条具有不同数据特征的代表序列，其余结果汇总于附录 E。

首先进行退化检验：令 $\phi=0$，fear-memory 的 ODE 右端项在 $x,y$ 分量上与 baseline 完全相同。数值验证（容差 $10^{-12}/10^{-14}$）得到两模型轨迹的最大偏差为 $1.48\times 10^{-3}$（相对变量尺度约 $0.007\%$）；该非零偏差来自两个系统维数不同导致的自适应步长差异，方程的 $x,y$ 动力学在结构上相同。

进一步分析各序列拟合所得 $\phi$ 与 baseline--fear-memory 的 $\Delta\mathrm{AICc}$ 关系。AICc 支持 baseline 的 13 条序列中，由于 6 条正式拟合选择零恐惧候选，fitted $\phi$ 的中位数约为 $6.0\times10^{-5}$；fear-memory 获得 AICc 支持的 2 条序列中，中位数约为 0.00196。当前样本未显示模型支持程度与 fitted $\phi$ 之间存在简单的单调关系，因此不能把较大的 fitted $\phi$ 解释为更强的机制证据。$\Delta\mathrm{AICc}$ vs fitted $\phi$ 散点图见附录 E。

$\phi$ 仅在 Andrén-2 的正式拟合中触及上界 0.2。自适应扩边诊断将上界提高到 1.0 后得到内部解 $\phi\approx0.766$，训练 RMSE 从 0.1011 降至 0.0981、AICc 从 $-156.5$ 降至 $-158.7$，但留后 RMSE 从 0.0833 上升至 0.0907。该结果说明原上界限制了训练拟合，却不能证明更大的 $\phi$ 具有更好的预测或生态解释；扩边结果仅作为边界敏感性诊断，不替换正式胜负。

固定 $M(0)=y(0)$ 和 $\delta=1$ 的理由如下：固定 $\delta=1$ 提供统一的记忆时间尺度，避免在短序列中与 $\phi$ 相互补偿；固定 $M(0)=y(0)$ 使初始记忆与首个可观测捕食压力一致，避免引入额外的弱可辨识初始状态。为检验上述设定是否主导结果，本文对全部 15 条序列进行 $M(0)/y(0)$ 与 $\delta$ 的 profile 分析（在每个固定值下重新优化其余参数，而非固定其余参数的条件扫描）。正文展示 Isle Royale、GLERL 与 TP：Isle Royale 为最长序列（38 年），GLERL 代表经前导零观测起点修正后仍呈宽记忆 profile 的浮游动物序列，TP 为 Killifish 研究内 fear-memory 留后验证优于 baseline 的代表。

$M(0)/y(0)$ profile 在 $[0,5]$ 上普遍平坦，全部 15 条序列的近似 95\% profile 集均同时触及扫描上下界，表明初始记忆的变化通常可由其他动力学参数补偿。触及边界表示可接受范围可能延伸到扫描区间之外，不应解释为精确置信区间端点。

\begin{figure}[H]
 \centering
 \begin{subfigure}{0.32\textwidth}
  \imginclude[width=\textwidth]{deep_analysis/tier1/m0_profile_isle_royale_wolf_moose_pre_2018.png}
  \caption{Isle Royale}
 \end{subfigure}\hfill
 \begin{subfigure}{0.32\textwidth}
  \imginclude[width=\textwidth]{deep_analysis/tier1/m0_profile_glerl_m110_zoop_1994-201.png}
  \caption{GLERL}
 \end{subfigure}\hfill
 \begin{subfigure}{0.32\textwidth}
  \imginclude[width=\textwidth]{deep_analysis/tier1/m0_profile_timeserieslogmeans_TP.png}
  \caption{TP}
 \end{subfigure}
 \caption{$M(0)/y(0)$ profile 的三条代表序列。蓝线为训练 RMSE，红线为留后 RMSE；宽而平坦的训练 profile 表明初始记忆弱可辨识。}
 \label{fig:m0-profile-body}
\end{figure}

$\delta$ 的 profile 显示，训练最优与留后最优的衰减时间尺度不同；14/15 条序列的近似 95\% profile 集至少触及一个扫描边界。在留后段上探索性选择最优 $\delta$，相对固定 $\delta=1$ 的验证 RMSE 改善中位数约为 1.2\%，其中 Windermere 北湖盆与 Isle Royale 的改善分别为 79.5\% 和 12.3\%，因此不能声称 $\delta$ 对所有序列均不重要。

\begin{figure}[H]
 \centering
 \begin{subfigure}{0.32\textwidth}
  \imginclude[width=\textwidth]{deep_analysis/tier1/delta_profile_isle_royale_wolf_moose_pre_2018.png}
  \caption{Isle Royale}
 \end{subfigure}\hfill
 \begin{subfigure}{0.32\textwidth}
  \imginclude[width=\textwidth]{deep_analysis/tier1/delta_profile_glerl_m110_zoop_1994-201.png}
  \caption{GLERL}
 \end{subfigure}\hfill
 \begin{subfigure}{0.32\textwidth}
  \imginclude[width=\textwidth]{deep_analysis/tier1/delta_profile_timeserieslogmeans_TP.png}
  \caption{TP}
 \end{subfigure}
 \caption{$\delta$ profile 的三条代表序列。训练最优与留后最优时间尺度可能不同，说明按训练拟合识别记忆尺度与按预测误差选择记忆尺度并不等价。}
 \label{fig:delta-profile-body}
\end{figure}

由于最优 $\delta$ 使用留后段选择，该分析只能作为敏感性评估，不构成新的无偏验证。

综上，$\phi$ 的边界触及以及全序列 $M(0)$、$\delta$ profile 的宽可接受范围共同表明，当前数据不足以支持对 fear-memory 恐惧强度和记忆尺度的稳定点估计。因此，本文关于机制差异的结论来源于\textbf{结构层面的模型间比较}（不同 ODE 右端项在相同数据下的相对表现），而非特定参数的点估计值。profile 图与完整 $\delta$ 比较表见附录 E。

\subsection{讨论：证据能够支持什么}
\label{sec:discussion}

本文的理论实验与数据拟合之间存在一种有意义的张力，这与文献中反复出现的方法论问题相呼应。Wang 等（2019）和 Zanette 等（2003）的工作表明恐惧效应在原理上可以改变种群动力学；Preisser 等（2007）和 Peacor \& Werner（2001）的系统综述确认了 NCE 的生态重要性。这些理论模型关心的是\textbf{可能存在什么}——它们证明，一旦恐惧以某种机制进入动力学方程，确实可以产生可观察的种群响应。然而，当参数被校准到真实数据时，问题转变为\textbf{在当前数据中实际被支持的是什么}。本文的六模型比较和全序列参数 profile 正是从这一角度切入：含恐惧模型可在部分序列上改善留后预测，但 AICc 在多数序列上支持较简单的 baseline，且 $M(0)$ 与 $\delta$ 的宽 profile 表明恐惧相关参数可被其他动力学参数补偿。因而，当前诊断不足以支持对真实恐惧强度的可靠估计。

数值实验支持”机制路径决定定性方向”，数据拟合支持”含恐惧模型能够改善部分序列的留后预测”，但二者都不足以证明真实系统存在相应恐惧机制。这一结论有三项独立的方法论依据。第一，15 条序列中 Andrén 七区、Killifish 三站和 Windermere 两湖盆存在研究内重复，统计独立性有限，不能视为 15 个独立生态证据。第二，AICc 在 baseline--fear-memory 两模型比较中有 13/15 条支持 baseline，在完整六模型比较中 baseline 有 8/15 条取得最低 AICc；留后预测优势与复杂度惩罚后的模型选择之间存在分歧。第三，记忆相关参数存在宽 profile，且 Windermere 案例揭示了外部人为干预（刺网移除）对封闭 ODE 假设的违反。综合这三点，预测能力与机制识别必须分开解释。

记忆参数的补充结果同样需要遵守这一证据边界。固定 $\delta=1$ 的主比较使用预先统一的时间尺度；附录中所谓”留后最优 $\delta$”是在同一留后段上比较多个候选值后得到，只能作为敏感性与假设生成分析，不能再视为无偏留后验证，也不用于重新计算模型胜负。

\medskip
\textbf{Windermere 案例：外部干预下的参数畸变。}Windermere 南北两湖盆的 baseline 与 fear-memory 留后验证 RMSE 分别为 1.52 和 3.04（北）与 1.16 和 1.39（南），在所有 15 条序列中偏高。两个湖盆的河鲈（猎物）在训练段内持续增长，而狗鱼（捕食者）处于刺网移除干预背景；该外部干预未进入本文的封闭 ODE。Windermere North 的 baseline 将 $e$ 压至下界 0.01，fear-memory 虽降低训练猎物 RMSE 至 0.049，却产生更差的留后外推。其他通道在该序列上的验证 RMSE 接近 0.423，但这种改善只能说明不同结构对未建模干预的适应能力不同，不能归因于某一恐惧机制更真实。

此案例的方法论含义超出了单条序列的拟合质量评估。当数据包含外部人为干预（刺网移除、补充投喂、栖息地管理等）时，封闭 ODE 模型会通过参数畸变来补偿未建模的过程。这意味着：(1) 极端拟合误差可能反映数据质量或模型假设违反，而非模型结构的生态合理性；(2) 畸变后的参数估计（如 $e\rightarrow 0$）不应被解读为真实生态速率；(3) 在此类序列中，不同恐惧通道之间的验证误差差异更应被理解为模型应对未建模扰动的灵活性差异，而非恐惧机制本身的优越性。因此，本文在主比较中保留 Windermere 序列以保持分析口径的统一性，但其拟合结果不独立用于支持或反对任何特定恐惧机制。

\section{结论}

本文围绕恐惧效应的机制路径依赖性、含恐惧模型的预测能力、以及预测改善能否等同于机制识别三个核心问题，得到以下分层结论。

\textbf{理论层面：恐惧效应存在机制路径依赖性。}在当前参考参数与 20\% 标准化基准下，五种 Holling II 恐惧机制均保持共存，但长期种群变化方向明显不同：繁殖抑制路径使捕食者均值下降 16\%--18\%，觅食抑制与处理时间延长则同时提高猎物和捕食者均值。换言之，种群是增加还是减少取决于恐惧通过何种机制进入动力学方程。该结论来自数值实验，证据强度为中等——20\% 是为清晰展示机制差异而人为选择的比较基准，不代表真实系统受恐惧抑制的实际强度，不同过程量之间也不能据此进行严格效应量排序。

\textbf{预测层面：含恐惧模型在部分序列上改善留后预测，但无单一通道一致胜出。}在 15 个拟合单位（序列级）的六模型统一比较中，留后验证最优的模型因序列而异；AICc 标准下 baseline 在 baseline--fear-memory 两模型比较中为 13/15，扩展至六模型后为 8/15。该结果不支持"添加恐惧一定改善种群预测"的普遍预期，也不支持"某一种恐惧机制普遍正确"。六个研究级单位的敏感性分析不改变这一结论。证据强度为中等——15 个拟合单位覆盖哺乳类、鱼类和浮游动物，但部分序列来自同一研究设计，统计独立性有限；同时，六模型拟合中 57/90 为 usable\_limit，胜负统计仍包含优化不确定性。

\textbf{机制层面：预测改善不能等同于恐惧机制的识别。}含恐惧模型在某些序列上留后预测更优，但三点限制因素决定了不能将此解释为"数据中存在恐惧效应"的正面证据。第一，AICc 在多数序列上支持不含恐惧的 baseline，与留后预测优势形成分歧；AICc 小样本修正只是在完整复杂度惩罚中增加额外校正，新增参数是否值得必须看完整 $\Delta\mathrm{AICc}$。第二，fear-memory 的 $\phi$ 在 Andrén-2 中触及正式上界，扩边虽改善训练拟合却恶化留后预测；全序列 $M(0)$ 与 $\delta$ profile 显示记忆相关参数存在宽可接受范围。固定 $M(0)$ 与 $\delta$ 是对弱可辨识性的控制，并不意味着它们对轨迹无影响。第三，Windermere 序列的案例表明，当数据包含外部人为干预（刺网移除）时，封闭 ODE 模型通过参数畸变补偿未建模过程。三种诊断从不同角度共同限定了本文能够提出的机制结论。

综合来看，本文的贡献不在于否定或证实恐惧效应的生态重要性，而在于：(1) 展示了五个恐惧通道在 20\% 标准化基准下的定性分歧，这意味着缺乏独立机制证据时，任何单一通道的拟合成功都不应解释为"找到了正确的恐惧机制"；(2) 在对 15 个公开时间序列的统一比较中，没有一个恐惧通道稳定优于不含恐惧的 baseline，说明此类宏观种群数据尚不足以一致区分恐惧通道；(3) 参数 profile 分析表明，即使恐惧模型在预测上有所改善，恐惧强度和记忆尺度仍可能存在宽可接受范围，因此相关点估计需要更审慎的解读。

\section*{参考文献}
\addcontentsline{toc}{section}{参考文献}
\begin{enumerate}[label={[\arabic*]}]
 \item Lotka A J. \emph{Elements of Physical Biology}. 1925.
 \item Rosenzweig M L, MacArthur R H. Graphical representation and stability conditions of predator--prey interactions. \emph{The American Naturalist}, 1963.
 \item Abrams P A. The consequences of adaptive behavior for population dynamics. \emph{Ecology}, 2000.
 \item Zanette L Y, et al. Perceived predation risk reduces offspring songbirds produce. \emph{PNAS}, 2003.
 \item Preisser E L, et al. Scared to death: intimidation and consumption in predator--prey interactions. \emph{Ecology}, 2007.
 \item Peacor S D, Werner E E. Trait-mediated indirect effects. \emph{Ecology}, 2001.
 \item Wang S, et al. Tipping points for predator--prey systems with fear effect. \emph{Ecology Letters}, 2019.
 \item Andrén H, et al. Eurasian lynx--roe deer multi-region dataset. Dryad, \url{https://doi.org/10.5061/dryad.9zw3r22mq}.
 \item Marino J A, et al. Lake Michigan zooplankton CE/NCE dataset. Dryad, \url{https://doi.org/10.5061/dryad.bh688ft}.
 \item Peacor S D, et al. Predation-risk effects literature review data. Dryad, \url{https://doi.org/10.5061/dryad.ffbg79cxt}.
\end{enumerate}

\clearpage
\appendix

\section{数据来源与预处理}

\subsection{主序列与分组}

\begin{longtable}{p{0.19\textwidth}p{0.14\textwidth}p{0.16\textwidth}p{0.38\textwidth}}
 \caption{15 条主序列清单与观测含义}\label{tab:datasets}\\
 \toprule
 简称 & 类群/分组 & 数据来源 & 猎物与捕食者观测\\
 \midrule
 \endfirsthead
 \toprule
 简称 & 类群/分组 & 数据来源 & 猎物与捕食者观测\\
 \midrule
 \endhead
 GLERL & zooplankton & Lake Michigan & \emph{Daphnia mendotae} 与 \emph{Bythotrephes}\\
 Andrén-1 至 7 & mammal & 同一 Andrén 研究 & 七个区域的狍猎获密度与猞猁家族群密度\\
 WRHW、TP、WRGP & fish & 同一 Killifish 研究 & 三站点 Least killifish 与 Eastern mosquitofish 月密度\\
 Isle Royale & mammal & NPS & 1980--2017 年驼鹿与狼数量；排除 2018 年迁入狼干预\\
 Windermere North/South & fish & GPDD & 两湖盆河鲈与狗鱼 kg/ha；狗鱼处于刺网移除干预背景\\
 Komi lynx--hare & mammal & GPDD & 1922--1942 年山兔与欧亚猞猁比例变换计数\\
 \bottomrule
\end{longtable}

Andrén 七区共享研究设计、物种和时间背景；Killifish 三站共享研究设计和物种；Windermere 两湖盆也来自同一研究。因此，15 条序列是 15 个拟合单位，而不是 15 个相互独立的生态重复。序列级胜负只描述当前样本中的模型表现。

项目另有 LTER 湖泊鱼类额外拟合，但其物种配对、时间尺度和筛选规则与主序列集合不同，仅作为探索性输出，不纳入本文 15 条主序列的胜负统计。

\subsection{预处理与质量控制}

数据加载器统一识别时间、猎物和捕食者列；Andrén 数据按面积转为密度，Killifish 对数均值还原到线性尺度。归一化尺度仅由训练段确定，避免验证信息泄漏。原始丰度 RMSE 仅用于解释单条序列，跨模型比较统一使用归一化 RMSE。数据文件解析还检查 Dryad 反爬页面并在必要时回退到项目内打包数据。

\section{数值方法与诊断}

\subsection{积分、长期收敛与灭绝}

自适应 RK45 解插值到默认 800 点均匀网格。长期均值和振幅取最终 20\% 窗口，并与前一相邻窗口比较；均值、振幅和漂移的相对变化均不超过 1\% 才标记为收敛。若未收敛，积分终点翻倍，最多三次。

灭绝阈值为对应系统尺度的 $10^{-3}$。只有尾段至少 80\% 低于阈值、终点仍低于阈值且最后恢复窗口没有越过阈值时，才判定相应种群灭绝。该规则避免把暂时低谷或不同量纲误判为灭绝。

\subsection{平衡点与残差校验}

Holling 默认平衡目标的诊断见附录 E。

\subsection{20\% 标准化恐惧校准}

20\% 抑制是人为选择的机制比较基准，用于使不同通道造成的动力学方向与幅度差异清晰可见。该数值不代表文献共识或真实系统中的典型恐惧强度；本节只比较在共同设定强度下，不同机制形式会产生何种结果。

\begin{table}[H]
 \centering
 \caption{参考平衡点处 20\% 标准化抑制参数}
 \label{tab:equiv-params}
 \small
 \begin{tabular}{lll}
  \toprule
  机制 & 匹配对象 & 校准参数\\
  \midrule
  fear-memory & Holling 增长项，$M^*=y^*/\delta$ & $\phi=0.00783431,\ \delta=1$\\
  fear-instant & Holling 增长项 & $\phi=0.00783431$\\
  fear-saturating & Holling 增长项 & $\phi=0.4,\ h=20.4172$\\
  fear-foraging & Holling 捕食率 & $\psi=0.0122446$\\
  fear-handling & Holling 捕食率 & $\psi=0.129848$\\
  \bottomrule
 \end{tabular}
\end{table}

Holling 无恐惧正平衡点为 $(20.0225,20.4172)$。数值校验中五种机制均达到 0.20 的目标抑制。

\section{完整机制与参数实验}

\subsection{机制比较}

\begin{table}[H]
 \centering
 \caption{20\% 标准化抑制基准下相对无恐惧基线的长期均值变化。memory reproduction 与 instant reproduction 的 $\phi$ 值相同（0.00783431），平衡点处恐惧项等价；两者差异仅来自记忆滞后。不同路径匹配不同过程量，数值幅度不作严格效应量排序。}
 \label{tab:mechanism-full}
 \begin{tabular}{lrrll}
  \toprule
  恐惧机制 & 猎物变化 & 捕食者变化 & 状态 & 收敛\\
  \midrule
  memory reproduction & $-0.0000\%$ & $-16.6667\%$ & 共存 & 是\\
  instant reproduction & $+0.0000\%$ & $-16.6667\%$ & 共存 & 是\\
  saturating reproduction & $+0.0000\%$ & $-18.0196\%$ & 共存 & 是\\
  foraging suppression & $+35.1503\%$ & $+23.2571\%$ & 共存 & 是\\
  handling extension & $+60.1151\%$ & $+36.0179\%$ & 共存 & 是\\
  \bottomrule
 \end{tabular}
\end{table}

主模型补充实验（情景对比、$\delta$ 扫描、五种机制时间序列）图见附录\,\ref{sec:theory-figures}。

\subsection{扫描摘要}

\begin{table}[H]
 \centering
 \caption{$\phi$ 与 $\delta$ 扫描完整摘要}
 \begin{tabular}{p{0.14\textwidth}p{0.18\textwidth}p{0.22\textwidth}p{0.34\textwidth}}
  \toprule
  扫描 & 网格 & 共存/收敛 & 结果摘要\\
  \midrule
  $\phi$ & $0$--$0.06$，31 点 & 31 未灭绝；20 收敛，11 未收敛 & 猎物均值约 20.02；捕食者均值从 20.42 降至 8.06\\
  $\delta$ & $0.2$--$4.0$，20 点 & 16 收敛，4 未收敛 & 猎物均值 20.02--24.06；捕食者均值 3.96--17.14\\
  \bottomrule
 \end{tabular}
\end{table}

\section{完整拟合结果}

\subsection{30 次拟合逐行结果}

表\,\ref{tab:fit-full} 列出全部 30 次拟合。两模型分别拟合 6、7 个动力学参数，全部结果均达到成功或达限可用状态。触边参数尤其是 $e$、$K$、$\phi$ 仍提示点估计存在边界敏感性。

\begin{landscape}
\scriptsize
\begin{longtable}{llrrrlll}
 \caption{15 条序列 $\times$ 2 模型拟合。训练与验证 RMSE 均为归一化总 RMSE；30 次拟合均可进入比较。}\label{tab:fit-full}\\
 \toprule
 序列 & 模型 & $\trainrmse$ & $\valrmse$ & AICc & 优化状态 & 可比较 & 触边参数\\
 \midrule
 \endfirsthead
 \toprule
 序列 & 模型 & $\trainrmse$ & $\valrmse$ & AICc & 优化状态 & 可比较 & 触边参数\\
 \midrule
 \endhead
 GLERL & base & 0.1708 & 0.2814 & -956.3 & 成功 & 是 & -- \\
 GLERL & memory & 0.1707 & 0.2815 & -954.4 & 成功 & 是 & -- \\
 Andrén-1 & base & 0.0601 & 0.0702 & -199.0 & 成功 & 是 & $e$ \\
 Andrén-1 & memory & 0.0596 & 0.0652 & -196.6 & 成功 & 是 & -- \\
 Andrén-2 & base & 0.1031 & 0.0968 & -157.9 & 成功 & 是 & $e$ \\
 Andrén-2 & memory & 0.1011 & 0.0833 & -156.5 & 达限可用 & 是 & $e,\phi$ \\
 Andrén-3 & base & 0.1130 & 0.0755 & -151.0 & 成功 & 是 & $K$ \\
 Andrén-3 & memory & 0.1131 & 0.0756 & -147.9 & 成功 & 是 & $K$ \\
 Andrén-4 & base & 0.0547 & 0.0652 & -194.3 & 达限可用 & 是 & $K$ \\
 Andrén-4 & memory & 0.0546 & 0.0653 & -191.3 & 达限可用 & 是 & $K$ \\
 Andrén-5 & base & 0.0618 & 0.0602 & -185.5 & 达限可用 & 是 & $K,e$ \\
 Andrén-5 & memory & 0.0618 & 0.0603 & -182.4 & 达限可用 & 是 & $K,e$ \\
 Andrén-6 & base & 0.0828 & 0.1810 & -144.0 & 达限可用 & 是 & $K$ \\
 Andrén-6 & memory & 0.0830 & 0.1812 & -140.6 & 达限可用 & 是 & $K$ \\
 Andrén-7 & base & 0.1057 & 0.2429 & -109.9 & 成功 & 是 & -- \\
 Andrén-7 & memory & 0.0753 & 0.1916 & -125.2 & 成功 & 是 & $e$ \\
 WRHW & base & 0.1835 & 0.1534 & -319.4 & 成功 & 是 & -- \\
 WRHW & memory & 0.1833 & 0.1536 & -317.3 & 成功 & 是 & -- \\
 TP & base & 0.1772 & 0.1030 & -319.3 & 达限可用 & 是 & $e$ \\
 TP & memory & 0.1767 & 0.0893 & -317.5 & 成功 & 是 & $e$ \\
 WRGP & base & 0.2052 & 0.1862 & -265.7 & 达限可用 & 是 & -- \\
 WRGP & memory & 0.2052 & 0.1862 & -263.3 & 达限可用 & 是 & -- \\
 Isle Royale & base & 0.1461 & 0.2283 & -217.3 & 达限可用 & 是 & -- \\
 Isle Royale & memory & 0.1461 & 0.2283 & -214.7 & 达限可用 & 是 & -- \\
 Windermere North & base & 0.1655 & 1.5234 & -114.6 & 达限可用 & 是 & $e$ \\
 Windermere North & memory & 0.1607 & 3.0392 & -113.6 & 达限可用 & 是 & -- \\
 Windermere South & base & 0.1405 & 1.1560 & -126.4 & 达限可用 & 是 & $a,e$ \\
 Windermere South & memory & 0.1389 & 1.3883 & -124.2 & 达限可用 & 是 & -- \\
 Komi lynx--hare & base & 0.3064 & 0.5811 & -60.3 & 成功 & 是 & -- \\
 Komi lynx--hare & memory & 0.2660 & 0.6085 & -66.1 & 达限可用 & 是 & $r,e$ \\
 \bottomrule
\end{longtable}
\end{landscape}

\subsection{验证与 AICc 胜负}

\begin{table}[H]
 \centering
 \caption{逐序列留后验证与 AICc 最优模型}
 \begin{tabular}{lll}
  \toprule
  序列 & 留后验证最优 & AICc 最优\\
  \midrule
  GLERL & baseline & baseline\\
  Andrén-1 & fear-memory & baseline\\
  Andrén-2 & fear-memory & baseline\\
 Andrén-3 & baseline & baseline\\
  Andrén-4 & baseline & baseline\\
 Andrén-5 & baseline & baseline\\
 Andrén-6 & baseline & baseline\\
  Andrén-7 & fear-memory & fear-memory\\
  WRHW & baseline & baseline\\
  TP & fear-memory & baseline\\
  WRGP & baseline & baseline\\
  Isle Royale & baseline & baseline\\
  Windermere North & baseline & baseline\\
  Windermere South & baseline & baseline\\
  Komi lynx--hare & baseline & fear-memory\\
  \midrule
 汇总 & base 9；memory 6；并列 0 & base 13；memory 2\\
  \bottomrule
 \end{tabular}
\end{table}

\subsection{六模型与研究级胜负}

\begin{table}[H]
 \centering
 \caption{完整六模型逐序列胜负。}
 \label{tab:six-model-winners}
 \small
 \begin{tabular}{lll}
  \toprule
  序列 & 留后验证最优 & AICc 最优\\
  \midrule
  GLERL & fear-handling & baseline\\
  Andrén-1 & fear-memory & fear-saturating\\
 Andrén-2 & fear-memory & fear-instant\\
 Andrén-3 & fear-foraging & baseline\\
 Andrén-4 & fear-foraging & fear-foraging\\
 Andrén-5 & baseline & baseline\\
 Andrén-6 & fear-handling & baseline\\
 Andrén-7 & fear-handling & fear-instant\\
  WRHW & fear-instant & fear-foraging\\
  TP & fear-handling & baseline\\
 WRGP & fear-memory & baseline\\
 Isle Royale & fear-memory & baseline\\
  Windermere North & fear-instant & fear-foraging\\
  Windermere South & baseline & baseline\\
  Komi lynx--hare & baseline & fear-memory\\
  \bottomrule
 \end{tabular}
\end{table}

六模型序列级留后赢家汇总为 baseline 3、fear-memory 4、fear-instant 2、fear-foraging 2、fear-handling 4；AICc 赢家汇总为 baseline 8、fear-memory 1、fear-instant 2、fear-saturating 1、fear-foraging 3。

\begin{table}[H]
 \centering
 \caption{研究内序列取中位数后的六个研究级单位胜负}
 \label{tab:study-level-winners}
 \small
 \begin{tabular}{p{0.22\textwidth}p{0.34\textwidth}p{0.25\textwidth}}
  \toprule
  研究级单位 & 留后验证最优 & AICc 最优\\
  \midrule
 Andrén & fear-foraging & fear-instant\\
  Killifish & fear-instant & baseline\\
  Windermere & fear-instant & baseline\\
  GLERL & fear-handling & baseline\\
 Isle Royale & fear-memory & baseline\\
  Komi lynx--hare & baseline & fear-memory\\
  \midrule
 汇总 & base 1；memory 1；instant 2；foraging 1；handling 1
         & base 4；memory 1；instant 1\\
  \bottomrule
 \end{tabular}
\end{table}

六模型 90 次拟合中，33 次为 success，57 次为 usable\_limit。零恐惧候选继承对应 baseline 的优化状态，因此该计数不会把未收敛的核心参数拟合重新标记为成功。仅保留 success 时，Windermere South 无模型可比较，其他序列保留的候选集合也不一致；因此 success-only 结果仅用于提示优化不确定性，不作为替代胜负表。

\subsection{不同模型均可能获胜的代表图}

\begin{figure}[H]
 \centering
 \begin{subfigure}{0.48\textwidth}
  \imginclude[width=\textwidth]{calibration_bda/figures/09_andren_lynx_roedeer_data_7_fear_memory.png}
  \caption{Andrén-7：memory 的 AICc 最优}
 \end{subfigure}\hfill
 \begin{subfigure}{0.48\textwidth}
  \imginclude[width=\textwidth]{calibration_bda/figures/14_windermere_south_pike_perch_baseline.png}
  \caption{Windermere South：baseline 验证最优}
 \end{subfigure}
 \caption{代表性拟合说明，不同模型可在不同序列和标准下获胜。}
\end{figure}

\section{参数可辨识性的补充限制}

\subsection{初始记忆与记忆衰减率的稳健性}

为检验固定 $M(0)=y(0)$ 与 $\delta=1$ 是否主导 fear-memory 的预测结果，本文进行两类补充分析。第一，在 $M(0)/y(0)\in[0,5]$ 上同时进行条件敏感性扫描和 profile：前者固定其余参数，后者在每个固定 $M(0)$ 下重新优化 $r,K,a,\theta,e,\mu,\phi$。第二，在 $\delta\in[0.01,10]$ 上固定每个 $\delta$ 并重新优化其余参数，同时报告对应记忆时间尺度 $1/\delta$、训练 profile 与连续多步留后 RMSE。

$M(0)$ 和 $\delta$ profile 分析的方法细节与代表图见正文第\,\ref{sec:identifiability} 节；此处保留 $\delta$ 比较的数值表格。表中固定 $\delta=1$ 的数值来自 profile 阶段的局部重新优化，不是正式多起点拟合结果的直接复现，因此个别序列与附录 D 的正式 fear-memory 指标略有差异。

\begin{landscape}
\begin{table}[H]
 \centering
 \caption{固定 $\delta=1$ 与探索性留后最优记忆尺度的比较}
 \label{tab:delta-holdout}
 \small
 \begin{tabular}{lrrrrr}
  \toprule
  序列 & $\valrmse(\delta=1)$ & 留后最优 $\delta$ & $1/\delta$ & 最优 $\valrmse$ & 改善比例\\
  \midrule
  GLERL & 0.2816 & 4.2170 & 0.24 & 0.2801 & 0.5\%\\
  Andrén-1 & 0.0652 & 0.2054 & 4.87 & 0.0580 & 11.0\%\\
  Andrén-2 & 0.0833 & 0.3162 & 3.16 & 0.0738 & 11.5\%\\
  Andrén-3 & 0.0748 & 1.0000 & 1.00 & 0.0748 & 0.0\%\\
  Andrén-4 & 0.0653 & 0.2054 & 4.87 & 0.0647 & 0.9\%\\
  Andrén-5 & 0.0610 & 0.1334 & 7.50 & 0.0603 & 1.0\%\\
  Andrén-6 & 0.1821 & 0.0154 & 64.94 & 0.1808 & 0.7\%\\
  Andrén-7 & 0.1912 & 1.0000 & 1.00 & 0.1912 & 0.0\%\\
  WRHW & 0.1536 & 0.2054 & 4.87 & 0.1504 & 2.1\%\\
  TP & 0.0893 & 4.2170 & 0.24 & 0.0832 & 6.8\%\\
  WRGP & 0.2050 & 0.2054 & 4.87 & 0.1876 & 8.5\%\\
  Isle Royale & 0.2288 & 0.0365 & 27.38 & 0.2006 & 12.3\%\\
  Windermere North & 3.1668 & 10.0000 & 0.10 & 0.6490 & 79.5\%\\
  Windermere South & 1.3890 & 0.3162 & 3.16 & 1.3788 & 0.7\%\\
  Komi lynx--hare & 0.6084 & 1.7783 & 0.56 & 0.6012 & 1.2\%\\
  \bottomrule
 \end{tabular}
\end{table}
\end{landscape}

在 15 条序列中，探索性选择其他 $\delta$ 后有 12 条留后 RMSE 数值下降，6 条改善至少 5\%，4 条改善至少 10\%，改善比例中位数为 1.2\%。因此，固定 $\delta=1$ 对总体模型比较结论影响有限，但少数序列可能包含不同记忆时间尺度的预测信号。由于最优 $\delta$ 使用留后段选择，表\,\ref{tab:delta-holdout} 不构成新的无偏验证，也不改变正文的模型胜负统计。

\subsection{Holling 默认目标与探索性证据}

Holling 默认平衡目标诊断中，15 条报告序列均未同时通过严格的稳定或联合周期判定，因此不存在合格的统一平衡目标。这并非数值失败：观测序列包含瞬态、环境噪声和代理指标，本就不应被强制解释为同一平衡点。即使构造目标，平衡关系也只能约束参数组合，存在结构不可辨识性。

Peacor 的 64 条记录仅为研究类型清单：TMIE 37、NCE 27。下图作为清单分布的补充展示，不解释为效应量分布。

\begin{figure}[H]
 \centering
 \imginclude[width=0.70\textwidth]{deep_analysis/tier2/peacor_effect_distribution.png}
 \caption{Peacor 分类字段的补充分布图；类别计数不是效应量。}
\end{figure}

珊瑚和豆娘结果仅记录为探索性数据。珊瑚位置变量与风险方向的映射不足，豆娘处理汇总中的活动方向也不能直接映射为繁殖抑制参数；二者均不进入定量结论。

\section{理论分析补充图}
\label{sec:theory-figures}

以下为正文理论与动力学分析及第\,\ref{sec:identifiability}\,节中引用的补充图表，移至此附录以保持正文聚焦核心证据。

\begin{figure}[H]
  \centering
  \imginclude[width=0.82\textwidth]{01_timeseries_baseline_vs_fear.png}
  \caption{Holling II 基线与恐惧记忆模型的时间序列对照。}
  \label{fig:main-ts}
\end{figure}

\begin{figure}[H]
  \centering
  \imginclude[width=0.80\textwidth]{05_phi_scan.png}
  \caption{$\phi$ 扫描：31/31 在当前积分范围内未灭绝，其中 11 个高 $\phi$ 情景未通过长期收敛诊断（不能视为已达平衡）。}
  \label{fig:phi-scan}
\end{figure}

\begin{figure}[H]
  \centering
  \imginclude[width=0.88\textwidth]{09_literature_mechanisms_bars.png}
  \caption{20\% 标准化抑制基准下，各恐惧机制相对无恐惧基线的长期变化。}
  \label{fig:mechanisms}
\end{figure}

\begin{figure}[H]
 \centering
 \begin{subfigure}{0.48\textwidth}
  \imginclude[width=\textwidth]{04_three_scenarios.png}
  \caption{三个主模型情景}
 \end{subfigure}\hfill
 \begin{subfigure}{0.48\textwidth}
  \imginclude[width=\textwidth]{06_delta_scan.png}
  \caption{$\delta$ 扫描}
 \end{subfigure}
 \caption{主模型补充实验。}
\end{figure}

\begin{figure}[H]
 \centering
 \imginclude[width=0.65\textwidth]{08_literature_mechanisms_timeseries.png}
 \caption{标准化基准下五种恐惧机制的时间序列。本节展示默认参数下的理论性质；跨路径数值幅度不作严格效应量排序。}
\end{figure}

\end{document}
